%%=============================================================================
%% Voorwoord
%%=============================================================================

\chapter*{\IfLanguageName{dutch}{Woord vooraf}{Preface}}%
\label{ch:voorwoord}

%% TODO:
%% Het voorwoord is het enige deel van de bachelorproef waar je vanuit je
%% eigen standpunt (``ik-vorm'') mag schrijven. Je kan hier bv. motiveren
%% waarom jij het onderwerp wil bespreken.
%% Vergeet ook niet te bedanken wie je geholpen/gesteund/... heeft

Het onderwerp van deze bachelorproef werd mij aangeboden door Arcelor Mittal Gent. Toen ik er voor het eerst van hoorde, wist ik meteen dat het een rijk en praktisch relevant thema was dat aansloot bij mijn interesse in mainframe-technologie. Tijdens de 14 weken dat ik aan deze migratie heb gewerkt, heb ik mij verdiept in CA~Gen-modellen, reverse-engineeringmethodes, en de implementatie van PL/I. Het project combineerde theoretische literatuurstudie met hands-on werk, het opstellen van vertaaltabellen, de analyse van action blocks en de uitvoering van een proof-of-concept — en daardoor heb ik zowel technische vaardigheden als methodologische inzichten opgedaan.

Deze migratie was niet mogelijk geweest zonder de deskundige begeleiding van mijn mentoren, Wouter Vermeerschen en Sabine Maes. Hun feedback op ontwerpkeuzes en code heeft de kwaliteit van dit werk aanzienlijk verhoogd. Mijn bijzondere dank gaat ook uit naar Didier Marichal, die de stageplaats mogelijk maakte en als co-promotor optrad, en naar Thomas Parmentier voor zijn waardevolle opmerkingen gedurende het hele traject. Verder wil ik de PLAPO-groep en collega’s bedanken voor hun praktische hulp en het delen van hun kennis.

Ik hoop dat dit onderzoek een concrete bijdrage levert aan het gestructureerd uitvoeren van CA~Gen→PL/I-migraties. Veel leesplezier.